\section{Evidência de IMP-ROLE-003 --- Etapa A (fundação M9 backend)}

\textbf{Feature:} \texttt{IMP-ROLE-003} --- autorização persistida no commit.

\textbf{Estado anterior:} \texttt{DIVERGENTE}. O backend autorizava por
\texttt{validarPermissao}, que decide somente pela claim JWT; a maioria das
mutações não relia atividade, versão e papel persistidos, e o escopo de domínio
era consultado fora da unidade transacional.

\textbf{Estado depois:} \texttt{DIVERGENTE}. Esta rodada entrega apenas a
fundação de autoridade persistida M9 e migra o primeiro conjunto restrito de
operações de usuários/papéis. As demais callables ainda não foram migradas, logo
a feature transversal permanece aberta e as contagens globais da matriz não
mudam.

\subsection{Contrato M9 derivado}

A decisão de autoridade base exige, de forma fail closed: autenticação;
\texttt{Usuarios/uid} existente e \texttt{ativo=true}; \texttt{versao\_permissoes}
persistida inteira e não negativa; claim \texttt{versao\_permissoes} presente e
igual à persistida; claim \texttt{roles} em lista; toda claim de papel
pertencente ao conjunto fechado e com concessão persistida cujo
\texttt{id\_usuario} coincide com o UID; e o papel exigido pela operação presente
entre os papéis afirmados e persistidos. Fontes: Seção~7/M9, a interpretação
normativa da Seção~10.2--10.4 (\textit{erratum} M9), CUE
\texttt{formal\_m9\_autorizacao} e Alloy M9
(\texttt{M9-INV-001/002/005/008}, \texttt{M9-WIT-014/017/018}) e a superfície já
fixada em \texttt{IMP-RULES-001}. Escopo específico de domínio (vínculo de
almoxarifado, ownership de professor) permanece fora da primitiva base.

\subsection{Primitiva implementada}

\begin{center}\small
\path{functions/src/auth.ts}
\end{center}

\begin{itemize}
  \item \texttt{PAPEIS\_CONHECIDOS}: conjunto fechado M9 como fonte única no
  backend, reutilizado por \path{functions/src/usuarios.ts}.
  \item \texttt{extrairClaimsAutoridade(request)}: valida apenas a estrutura das
  claims (auth, lista de papéis, conjunto fechado, versão inteira não negativa).
  \item \texttt{resolverAutoridadePersistidaTx(tx, claims, papeisRequeridos)}:
  núcleo transacional; lê, na mesma transação do efeito, \texttt{Usuarios/uid} e
  cada documento de papel, e decide. Retorna
  \texttt{\{uid, papeis, versaoPermissoes, papelAutorizado\}}.
  \item \texttt{validarAutoridadePersistida}/
  \texttt{validarAutoridadePersistidaComClaims}: conveniências para pré-checagem
  antes de efeitos externos, reexecutando a leitura numa transação.
  \item \texttt{validarPermissao} permanece marcada como \textbf{LEGACY AUTH
  PATH}, explicitamente insuficiente para mutações críticas M9.
\end{itemize}

\subsection{Casos permanentes e ciclo RED/GREEN}

\begin{center}\small
\path{functions/src/__tests__/integration/m9-backend-authority.integration.test.ts}
\end{center}

Os casos usam Firestore Emulator e Admin SDK (não são mocks), com claims e
\texttt{request} construídos determinísticamente. A tabela registra a fonte
normativa, o comportamento do mecanismo anterior (RED) e o resultado após a
implementação (GREEN).

\small
\begin{longtable}{@{}p{0.18\textwidth}p{0.38\textwidth}p{0.13\textwidth}p{0.11\textwidth}@{}}
\toprule
ID & Obrigação observada & Antes & Depois \\
\midrule
\endhead
TEST-INT-M9-AUTH-001 & Não autenticado nega & PASS controle & PASS \\
TEST-INT-M9-AUTH-002 & UID sem \texttt{Usuarios} nega & FAIL & PASS \\
TEST-INT-M9-AUTH-003 & Usuário inativo nega & FAIL & PASS \\
TEST-INT-M9-AUTH-004 & Claim sem versão nega & PASS controle & PASS \\
TEST-INT-M9-AUTH-005 & Versão obsoleta nega & FAIL & PASS \\
TEST-INT-M9-AUTH-006 & Versão corrente + papel persistido autoriza & PASS controle & PASS \\
TEST-INT-M9-AUTH-007 & Papel permitido ausente nega & PASS controle & PASS \\
TEST-INT-M9-AUTH-008 & Papel desconhecido nega & PASS controle & PASS \\
TEST-INT-M9-AUTH-009 & Claim de papel sem concessão persistida nega & FAIL & PASS \\
TEST-INT-M9-AUTH-010 & Papel persistido não afirmado na claim nega & PASS controle & PASS \\
TEST-INT-M9-AUTH-011 & Documento de papel com UID incompatível nega & FAIL & PASS \\
TEST-INT-M9-AUTH-012 & \texttt{roles} não-array nega & PASS controle & PASS \\
TEST-INT-M9-AUTH-013 & Multi-role válido (Aluno+Bolsista) autoriza & PASS controle & PASS \\
TEST-INT-M9-AUTH-014 & Multi-role com papel desconhecido nega & PASS controle & PASS \\
\bottomrule
\end{longtable}
\normalsize

O RED foi executado antes da implementação, contra o mecanismo anterior: 6 de 19
casos falharam (002, 003, 005, 009, 011 e OP-002), evidenciando precisamente a
aceitação de token sem atividade, sem versão corrente e sem papel persistido.
Após a implementação, os 19 casos passam.

\subsection{Operação real protegida e ausência de escrita parcial}

Além da primitiva, \path{functions/src/usuarios.ts} foi migrado para a nova
autoridade transacional nas operações \texttt{convidarUsuario},
\texttt{revogarUsuarioPapel} e no núcleo \texttt{alterarPapel}. O bloco seguinte
do mesmo arquivo de teste prova a operação real.

\small
\begin{longtable}{@{}p{0.24\textwidth}p{0.58\textwidth}@{}}
\toprule
ID & Resultado \\
\midrule
\endhead
TEST-INT-M9-AUTH-OP-001 & Chefia corrente e persistida autoriza; revogação aplicada. \\
TEST-INT-M9-AUTH-OP-002 & Token obsoleto nega e o Firestore permanece inalterado. \\
TEST-INT-M9-AUTH-OP-003 & Chefia inativa nega e o Firestore permanece inalterado. \\
TEST-INT-M9-AUTH-OP-004 & Claim de Chefe sem documento persistido nega e não muta. \\
\bottomrule
\end{longtable}
\normalsize

Cada caso negado verifica o estado persistido posterior, não apenas o
\texttt{HttpsError}: o documento do papel do alvo e o \texttt{ativo} permanecem.

\subsection{Ambiente e estratégia Node 20}

O backend declara \texttt{engines.node=20}. A evidência final foi executada em
Node \textbf{20.19.6} (npm 10.8.2). O \texttt{nixpkgs} atual removeu
\texttt{nodejs\_20} por EOL em 30/04/2026; a execução usou
\path{nix shell github:nixos/nixpkgs/nixos-25.05#nodejs_20} como fornecimento
temporário, sem alterar \texttt{flake.lock} e sem alterar o Home Manager:

\begin{center}
\texttt{NODE20\_STRATEGY = TEMPORARY\_NIX\_ONLY}; \quad
\texttt{HOME\_MANAGER\_CHANGED = NO}
\end{center}

\subsection{Resultados e regressão}

\small
\begin{longtable}{@{}p{0.52\textwidth}p{0.38\textwidth}@{}}
\toprule
Verificação (Node 20.19.6) & Resultado \\
\midrule
\endhead
\texttt{npm run build} & PASS \\
\texttt{npm run test:unit} & PASS (53/53) \\
\texttt{TEST-INT-M9-AUTH-*} + \texttt{-OP-*} (emulator) & PASS (19/19) \\
\texttt{usuarios.integration.test.ts} (fixtures M9) & PASS (2/2) \\
\texttt{npm run test:integration} & 90 passam / 14 falham \\
\bottomrule
\end{longtable}
\normalsize

A baseline registrada era 69 passam / 16 falham. A diferença é integralmente
explicada: os 19 casos M9 são novos e passam, e 2 casos legados de
\texttt{usuarios.integration.test.ts} que já falhavam por fixture sem autoridade
persistida passaram a ter arrange M9 válido. As 14 falhas remanescentes estão
somente em \texttt{reagentes} (10), \texttt{patrimonio} (2) e \texttt{roteiros}
(2), nenhuma em \texttt{auth}/\texttt{usuarios}; nenhuma falha nova não explicada.

O lint direcionado aos arquivos alterados não introduz novos erros; os avisos
remanescentes (\texttt{no-explicit-any} em \texttt{catch} legado e no teste
legado) já existiam.

\subsection{Escopo e pendências}

Migrado nesta etapa: \texttt{alterarPapel}, \texttt{convidarUsuario} e
\texttt{revogarUsuarioPapel}, com a decisão de chefia relida na transação do
efeito (fecha TOCTOU). A ordem de \texttt{convidarUsuario} mantém efeitos de
Firebase Auth após uma pré-checagem M9 e antes da revalidação transacional; o
risco residual (criação de identidade Auth entre pré-checagem e commit) fica
registrado como dependência futura, sem redesenho do workflow de convite nesta
rodada.

Não migrado: reagentes, patrimônio, turmas, posts, roteiros, notificações,
relatórios, etiquetas; e \texttt{IMP-ROLE-001}, \texttt{IMP-ROLE-002} e
\texttt{IMP-IDEM-001} permanecem intocados. \texttt{IMP-ROLE-003} continua
\texttt{DIVERGENTE} por desenho.
